% !TITLE 诗经·小雅·采薇（节选）
% !AUTHOR 无名氏
% !DYNASTY 先秦
% !GENRE 四言诗
% !ORDER 4

\begin{poem}
昔我往矣，杨柳依依\textsuperscript{1}。
今我来思\textsuperscript{2}，雨雪霏霏\textsuperscript{3}。
行道迟迟\textsuperscript{4}，载渴载饥\textsuperscript{5}。
我心伤悲，莫知我哀！
\end{poem}

\begin{annotations}
\notetext{1}{依依：柔弱摇曳的样子。形容柳条随风轻拂，也暗喻离别时不舍的心情。}
\notetext{2}{思：句末助词，无实际意义，只起协调音节的作用。}
\notetext{3}{雨雪霏霏：雨，这里作动词，指下雪。霏霏，雪花大而密的样子。}
\notetext{4}{迟迟：迟缓、缓慢的样子。形容路途泥泞难行，也暗示战士近乡情怯的沉重心情。}
\notetext{5}{载：又、且。载渴载饥形容战士归途上又饥又渴。}
\end{annotations}

\begin{appreciation}
《采薇》全诗六章，写一个戍边战士从出征到归乡的全过程。这里节选的是末章，写他解甲归田、踏上归途的情景。

“昔我往矣，杨柳依依。今我来思，雨雪霏霏。”四句被东晋名将谢玄（草木皆兵、风声鹤唳的故事里的那位东晋大将军）认为是《诗经》里最美的句子。战士在美好的春天里离开家乡踏上战场，柳条温柔地随风摇曳好似对离人的挽留。“柳”与“留”同音，柳絮又纷纷似雪，因此古人常在离别时用柳树表达愁绪，后世也有折柳送别的习俗。在这样美好的季节里，在万物生发的春天中，战士却要踏上残酷的九死一生的战场，春天的美好反衬出了战士的哀愁，这就是以乐写哀。

战争非一日之功。战士春天上战场，直到漫天飞雪的冬天才回来。经历战火而身心俱疲的战士在肃杀的冬天回来，环境的凄凉也映衬出战士内心的悲伤与疲惫。于是往下便是直抒胸臆，明写战士内心的疲惫与哀愁：

他在路上走得很慢，又渴又饿。这不仅是身体的疲惫，更是心理的沉重。他在拖延什么呢？也许是近乡情怯，也许是不知道家中亲人是否还在。"我心伤悲，莫知我哀"，他的悲哀，没有人真正理解。那些没有经历过战争的人，怎么可能体会到一个战士内心深处的创伤？

战争之下没有赢家。今天的世界也并不和平。贪婪的野心家们也许还会一次又一次地发动战争，采薇的哀歌也许也还会一直唱下去，但如果多一个人能理解这跨越了三千年的悲哀，或许和平的到来就会快一步。
\end{appreciation}
